\section{Introduction}

Foundation-model applications are no longer single models queried in isolation. A deployed assistant, analyst tool, or enterprise chatbot may combine retrieval-augmented generation (RAG), dense and lexical indexes, rerankers, prompt templates, supervised fine-tuning, LoRA adapters, synthetic-data generation, feedback traces, caches, monitoring logs, and human-facing citations. RAG has become a standard way to connect language models to external knowledge sources~\cite{lewis2020rag,karpukhin2020dpr,izacard2021fid,gao2024ragsurvey}, while instruction tuning and adapter-based specialization make private or domain-specific records part of model behavior. This architecture improves utility, but it also makes deletion requests technically ambiguous: deleting a raw document, account record, or support transcript does not necessarily remove the chunks, embeddings, summaries, synthetic question-answer pairs, cache entries, adapters, or evaluation traces that were derived from it.

This ambiguity is becoming a central data-management problem for foundation-model systems. Privacy and governance regimes motivate data erasure and deletion accountability~\cite{voigt2017gdpr,cohen2022deletioncontrol}, and recent foundation-model challenge work identifies private data use, continual learning, unlearning, watermarking, and efficiency as important open problems~\cite{fan2025fedfmchallenges}. Machine unlearning provides the algorithmic foundation for making systems forget selected data~\cite{cao2015machineunlearning,ginart2019making,bourtoule2021machineunlearning,guo2020certified}. However, in modern LLM and RAG deployments, the deletion problem is not confined to the parameter vector. A post-deletion system can pass a model-centric unlearning test while still retrieving a stale chunk, citing a shadow copy, answering from a synthetic derivative, or paraphrasing a fact learned through an adapter.

The difficulty is that deletion is a behavioral property of a pipeline, not merely a storage operation or a weight update. Existing LLM-unlearning benchmarks and methods increasingly study forgetting and retained utility for adapted language models~\cite{maini2024tofu,yao2024largeunlearning,liu2024rethinking,wang2024llmunlearning,jin2024surveyunlearning}. Separately, extraction and membership-inference work shows that language models can reveal training data or expose membership signals through black-box interaction~\cite{carlini2019secretsharer,carlini2021extracting,carlini2023quantifying,shokri2017membership,yeom2018privacy,mattern2023membership}. RAG evaluation studies retrieval quality, attribution, grounding, and context usage~\cite{gao2023retrievalaugmented,asai2024selfrag,yu2024ragevalsurvey,akkiraju2024facts}. These lines of work are all necessary, but none directly answers the operator's audit question: after a deletion request, what measurable evidence shows that the deleted data and its downstream derivatives no longer influence retrieval, generation, or fine-tuned behavior?

We argue that deletion verification should be treated as an auditable data-centric systems problem. The relevant object is not only a model checkpoint; it is the collection of data-bearing artifacts that may preserve a deleted record's influence. These artifacts include retriever documents, chunk boundaries, embedding rows, reranker candidates, prompt caches, summaries, synthetic examples, adapter updates, canaries, provenance tokens, and generated traces. A practical audit must therefore combine three ingredients. First, it needs provenance: which artifacts are downstream of the deletion target? Second, it needs prioritization: which deleted records had high influence or many derivatives and therefore deserve more audit budget? Third, it needs behavioral tests: does the post-deletion system still reproduce, retrieve, paraphrase, infer, or cite the deleted information?

This paper proposes \method{} (\emph{Auditable Data-Centric Unlearning}), a framework for deletion-risk auditing in retrieval-augmented and fine-tuned foundation-model systems. \method{} represents deletion targets and their downstream artifacts in a provenance graph; assigns audit priority using pre-deletion influence and provenance degree; generates direct, paraphrase, retrieval, counterfactual, extraction, trigger, and graph-pivot probes; and scores observed behavior across multiple residual-dependence channels. Instead of issuing a binary claim of perfect forgetting, \method{} returns a confidence-bounded pass, fail, or escalate decision under an operator-chosen tolerance. The goal is not to replace unlearning algorithms, certified removal mechanisms, differential privacy, or legal review~\cite{dwork2006dp,guo2020certified}. The goal is to provide measurable technical evidence about residual behavioral dependence in the data pipeline surrounding those mechanisms.

The core deletion-risk metric aggregates six evidence channels: direct leakage, paraphrased leakage, retrieval dependence, counterfactual answer dependence, extraction risk, and watermark or provenance-token hits. This design is motivated by known failure modes. Exact-match audits can miss paraphrased leakage; canary-only audits can fail when obvious markers are suppressed; membership tests can miss retrieval-only residues; and retriever-only audits can miss model-side memory or answer-side paraphrases. Recent work on RAG security and extraction further shows that leakage can arise through benign-looking queries, triggered behavior, and graph or retrieval pivots rather than only through direct requests for secrets~\cite{wang2025ikea,xie2026canaryrag,song2026grasp,thornton2026retrievalpivot}. \method{} therefore treats deletion verification as a multi-channel audit rather than a single memorization test.

We instantiate the framework in \bench{}, a reusable benchmark and artifact for deletion auditing. The benchmark covers four complementary settings. \emph{RealRAG} uses public FEVER evidence as a retain corpus while injecting synthetic private deletion targets and downstream derivatives. \emph{NaturalFEVER} treats public FEVER evidence passages themselves as deletion units, producing a natural-data deletion case with redacted outputs. \emph{FineTuneSim} provides a controlled adapter-memory simulator for repeated unlearning stress tests. \emph{HybridReal} combines RAG records and derivative artifacts that survive through summaries, caches, synthetic examples, or model-side behavior. We further add LoRA-SFT, cached PEFT/LoRA validations, pretrained LoRA margin analysis, TOFU-style public forget/retain evaluation, neural dense retrieval with MiniLM/E5/BGE-style retrievers, and a local cross-encoder-style reranking study.

Across the repeated submission-scale harness, \method{} detects configured deletion failures in RealRAG, NaturalFEVER, and HybridReal while preserving zero false alarms in the configured clean cases; it detects 97.2\% of FineTuneSim failures where narrow baselines miss important channels. The advanced experiments show the same pattern in model-side and retrieval-side settings: retriever-only audits do not detect LoRA memory, exact matching is brittle under paraphrase, and canary-only tests fail when semantic facts remain without the marker. The point of these experiments is not that \method{} proves legal compliance or absolute absence of influence. Rather, the results show that deletion audits become substantially more informative when they are provenance-aware, valuation-guided, and multi-channel.

This paper makes the following contributions:
\begin{itemize}
\item We formulate deletion verification for RAG and fine-tuned foundation-model systems as a pipeline-level data audit over provenance-linked artifacts, rather than as only model-weight unlearning.
\item We define a deletion-risk metric that combines direct leakage, paraphrased leakage, retrieval dependence, counterfactual answer dependence, extraction risk, and watermark or provenance-token evidence.
\item We present a valuation-guided black-box audit algorithm that allocates limited probe budgets toward high-risk deletion targets and returns confidence-bounded pass, fail, or escalate decisions.
\item We introduce \bench{}, a reusable benchmark spanning synthetic private deletion, natural public-evidence deletion, RAG retention, fine-tuning residues, hybrid derivative retention, dense retrieval, pretrained LoRA, and TOFU-style public forget/retain evaluation.
\item We provide an attack and baseline suite showing where exact-match, canary-only, retriever-only, membership-inference, extraction-only, and influence-proxy audits fail, together with redacted case studies and artifact code for reproduction.
\end{itemize}

The remainder of the paper is organized as follows. Section~2 discusses related work in machine unlearning, RAG, data valuation, extraction, watermarking, and deletion governance. Section~3 defines the audit problem and threat model. Sections~4 and~5 present the deletion-risk metric and audit algorithm. Section~6 describes \bench{}. Section~7 reports the experiments, baselines, ablations, and case studies. Sections~8 and~9 discuss limitations, artifact release, and reproducibility, and Section~10 concludes.
